
Blok şifrelerin kriptoanalizi, simetrik kriptografi alanında hâlen en kritik araştırma konularından biri olarak önemini korumaktadır. Özellikle son yıllarda gelişen gömülü sistem teknolojileri, Nesnelerin İnterneti (IoT) uygulamaları ve RFID sistemleri gibi kaynak kısıtlamalarının baskın olduğu alanlarda kriptografik algoritmaların etkinliği ve verimliliği daha da ön plana çıkmıştır. Bu tür ortamlarda kullanılan cihazlar, genellikle sınırlı işlem gücüne, düşük enerji kapasitesine ve kısıtlı bellek kaynaklarına sahiptir. Bu nedenle, standart şifreleme algoritmaları — örneğin AES — bu sistemlerde kullanılmak istendiğinde performans sorunları veya aşırı kaynak tüketimiyle karşılaşılmaktadır. Bu gerçeklik, hafif blok şifreleme algoritmalarının geliştirilmesini ve değerlendirilmesini zorunlu kılmaktadır.

Hafif şifreleme alanında öne çıkan alternatif algoritmalardan biri ITUbee blok şifre algoritmasıdır. ITUbee, hem yapısal sadeliği hem de düşük kaynak tüketimi ile özellikle kaynak kısıtlı ortamlarda kullanılmak üzere tasarlanmıştır. Ancak yeni bir şifreleme algoritmasının geliştirilmesi kadar, onun güvenliğinin sistematik ve derinlemesine incelenmesi de büyük önem arz etmektedir. Kriptografi literatüründe kabul gören standartlara göre, bir algoritmanın etkinliği yalnızca performans ölçütleriyle değil, aynı zamanda güçlü kriptanalitik dayanıklılık ile de değerlendirilmektedir.

Bu noktada, literatürde en yaygın ve etkili kriptoanaliz yöntemleri olan diferansiyel kriptoanaliz ve doğrusal kriptoanaliz öne çıkmaktadır. Bu saldırı teknikleri, blok şifrelerin yapısal zayıflıklarını tespit ederek, algoritmanın rastgeleliğe ne ölçüde yaklaştığını analiz etmeyi amaçlar. Diferansiyel kriptoanaliz, belirli giriş farklarının çıkışlara nasıl yansıdığını inceleyerek, yüksek olasılıklı fark yolları (differential trails) üzerinden anahtar bilgisi elde etmeye çalışırken; doğrusal kriptoanaliz, giriş ve çıkış bitleri arasında doğrusal bağıntılar kurarak benzer bir bilgi sızdırma tekniği uygular. Her iki yöntemin de başarılı olabilmesi için, şifreleme algoritmasında belirli yapısal düzenliliklerin veya zayıflıkların mevcut olması gerekir.

Bu iki yönteme ek olarak, ilişkili anahtar diferansiyel saldırılar (related-key differential attacks) da  etkin kriptoanaliz yöntemleri arasındadır. Bu saldırı türü, saldırganın belirli farklara sahip birden fazla anahtar üzerinden işlem yapabilmesine olanak tanıyan genişletilmiş bir modeldir. Özellikle basit veya yapılandırılabilir anahtar genişletme mekanizmalarına sahip algoritmalar, bu tür saldırılara karşı savunmasız kalabilmektedir.

Ancak, modern blok şifre algoritmalarının karmaşık yapıları göz önünde bulundurulduğunda, manuel (el ile) analiz yöntemlerinin uygulanabilirliği büyük ölçüde azalmaktadır. S-box yapısında veya doğrusal dönüşüm katmanlarında yapılan küçük bir değişiklik bile tüm analiz sürecinin baştan tasarlanmasını gerektirir. Bu durum, kriptoanaliz süreçlerinin otomatikleştirilmesini ve daha verimli bir biçimde yürütülmesini gerekli kılmaktadır. Bu bağlamda, Karma Tamsayılı Doğrusal Programlama (KTPL) temelli kriptanaliz yaklaşımların önemi artmıştır.

KTPL, doğrusal programlamanın klasik yapısını tamsayı değişkenlerle birleştirerek daha karmaşık karar problemlerini çözebilen güçlü bir optimizasyon tekniğidir. Kriptoanaliz bağlamında KTPL, bir şifreleme algoritmasının bileşenlerini (S-box, XOR, doğrusal dönüşüm, taşıyıcı matrisler vb.) doğrusal eşitsizlikler aracılığıyla matematiksel olarak modelleme imkânı sunar. Bu sayede, kriptoanalitik saldırılar bir optimizasyon problemi haline getirilir ve mevcut yazılımlar aracılığıyla etkili biçimde çözüm aranabilir. Özellikle hafif blok şifrelerde, KTPL modelleri düşük tur sayıları için tüm olası fark yollarını tarayabilecek kapasitededir. Bu da KTPL'yi, klasik arama yöntemlerine (örneğin brute-force, ağaç arama, rastgele örnekleme) kıyasla daha etkili ve sistematik bir araç haline getirmektedir.

Tez kapsamında, öncelikle simetrik şifreleme algoritmalarının temel yapı taşları açıklanmıştır. Substitution-Permutation Network (SPN) yapıları, Feistel ağları, ARX tabanlı şifreler ve sünger yapılar (sponge constructions) gibi yaygın kullanılan mimariler ele alınmış; bu yapıların güvenlik ve performans üzerindeki etkileri ayrıntılı bir biçimde incelenmiştir. Ardından, doğrusal, diferansiyel ve ilişkili anahtar diferansiyel kriptoanaliz teknikleri teorik temelleriyle birlikte değerlendirilmiştir. Bu bağlamda, kriptografik saldırıların temel varsayımları, başarı kriterleri ve uygulama koşulları detaylandırılmıştır.

Devamında, doğrusal programlama (LP) ve KTPL hakkında genel bilgiler sunulmuş; ardından KTPL’nin kriptoanaliz bağlamındaki kullanımı üzerinde durulmuştur. Özellikle “dallanma sayısı” (branch number) ve “H-temsili” (H-representation) gibi kavramlar tanıtılmış, bu yapıların doğrusal eşitsizliklerle nasıl modellenebileceği örnekler ile açıklanmıştır. Branş sayısı, doğrusal dönüşüm matrislerinin difüzyon kapasitesini ölçen kritik bir metrik olup, güvenlik değerlendirmelerinde sıklıkla kullanılmaktadır. H-temsili ise doğrusal olmayan bileşenlerin (örneğin S-box’ların) doğrusal kısıtlarla ifade edilmesine olanak tanıyan geometrik bir modelleme yaklaşımıdır.

Tezde ayrıca, KTPL modellerinde kullanılacak eşitsizlik sayısını azaltmaya yönelik olarak geliştirilen iki algoritma — ''Greedy'' ve ''New-Reduction'' — tanıtılmıştır. Bu algoritmalar, KTPL modellemelerinde performansı artırmak ve çözüm sürelerini düşürmek için kritik öneme sahiptir. Her bir algoritma, farklı hedef fonksiyonlara göre çalışmakta olup, çözüm uzayında daha verimli arama yapmayı mümkün kılmaktadır.

Tüm bu teorik temeller ve araçlar kullanılarak, tezde ITUbee blok şifre algoritması detaylı bir biçimde modellenmiş ve analiz edilmiştir. ITUbee, Feistel yapısına dayalı, 64 bit blok boyutuna sahip ve 128 bit anahtar kullanan bir algoritmadır. Tasarımı itibarıyla hafif cihazlar için optimize edilmiştir. Tezde, ITUbee’nin her bir bileşeni — S-box, doğrusal katman, anahtar genişletme mekanizması, Feistel yapısı vb. — KTPL modelleri ile ifade edilmiş ve bu modeller üzerinden saldırı karakteristikleri elde edilmiştir.

Elde edilen sonuçlara göre, ITUbee algoritması üç çevrimden sonra diferansiyel ve doğrusal saldırılara karşı güvenlidir. Bu, algoritmanın yazarlarının teorik olarak sunduğu güvenlik varsayımlarını desteklemektedir. Ancak ilişkili anahtar diferansiyel saldırıları açısından yapılan analizler, ITUbee algoritmasının yazarlarının teorik olarak gösterdiği 10 çevrimden daha güvenli olduğu, 8 çevrim ve sonrası için bu saldırı yönteminin uygulanamayacağı göstermiştir. 

Elde edilen karakteristikler, tezde tablo ve grafiklerle detaylı bir biçimde sunulmuş; kullanılan KTPL modelleri ise ekler kısmında verilmiştir. Bu çalışma, yalnızca ITUbee algoritmasının güvenliğini değerlendirmekle kalmayıp, aynı zamanda KTPL temelli modelleme yöntemlerinin pratikte nasıl uygulanabileceğini gösteren sistematik bir rehber sunmaktadır. Ayrıca, KTPL'nin hafif blok şifrelerin analizi için güçlü ve ölçeklenebilir bir araç olduğunu ortaya koymakta, gelecekteki kriptografik analiz çalışmaları için önemli bir zemin hazırlamaktadır.

 